\documentclass[12pt,a4paper]{article}

\usepackage[margin=1in]{geometry}
\usepackage{graphicx}
\usepackage{float}
\usepackage{caption}
\usepackage{xcolor}
\usepackage{listings}
\usepackage[hidelinks]{hyperref}

\graphicspath{{images/}}

\lstset{
    basicstyle=\ttfamily\small,
    breaklines=true,
    frame=single,
    numbers=left,
    numberstyle=\tiny,
    keywordstyle=\color{blue},
    commentstyle=\color{gray},
    stringstyle=\color{teal}
}

\newcommand{\screenshot}[3]{
    \begin{figure}[H]
        \centering
        \IfFileExists{images/#1}{
            \includegraphics[width=0.95\textwidth]{#1}
        }{
            \fbox{
                \begin{minipage}[c][0.18\textheight][c]{0.9\textwidth}
                    \centering
                    Missing screenshot file: \texttt{images/#1}
                \end{minipage}
            }
        }
        \caption{#2}
        \label{#3}
    \end{figure}
}

\begin{document}

\begin{titlepage}
    \centering
    \vspace*{2cm}
    {\Large \textbf{Student Course Registration System}}\\[0.4cm]
    {\large Spring MVC Data Binding and Response Experiment}\\[1.5cm]
    {\large Academic Laboratory Report}\\[2cm]
    \vfill
    \begin{flushleft}
        \textbf{Course:} Java EE / Spring MVC\\
        \textbf{Experiment Topic:} Data Binding and Response Handling\\
        \textbf{Project Name:} Data\_Binding\_and\_Response\\
        \textbf{Server:} Jetty 9.4.54\\
    \end{flushleft}
    \vfill
    {\large \today}
\end{titlepage}

\tableofcontents
\newpage

\section{Introduction}

This experiment implements a simple Student Course Registration System using Spring MVC.
The main purpose of the experiment is to understand how request data is received by controller
methods and how Spring MVC returns page and JSON responses. The system allows a student to submit
basic information, select multiple courses, and receive JSON data from the server.

\section{Experiment Objectives}

The objectives of this experiment are:

\begin{itemize}
    \item To practice simple parameter binding in Spring MVC.
    \item To bind request parameters using \texttt{@RequestParam}.
    \item To receive REST-style values using \texttt{@PathVariable}.
    \item To bind request parameters to a Java POJO.
    \item To bind multiple request values using a list.
    \item To receive JSON request data using \texttt{@RequestBody}.
    \item To return JSON response data using \texttt{@ResponseBody}.
\end{itemize}

\section{Development Environment}

\begin{table}[H]
    \centering
    \begin{tabular}{|l|l|}
        \hline
        \textbf{Item} & \textbf{Description} \\
        \hline
        Programming Language & Java \\
        Framework & Spring MVC \\
        Build Tool & Maven \\
        Server & Jetty 9.4.54 \\
        Testing Tool & Browser and Postman \\
        Project Type & Maven WAR Web Application \\
        \hline
    \end{tabular}
    \caption{Development environment}
\end{table}

\section{System Design}

The application uses the Spring MVC pattern. The \texttt{StudentController} receives HTTP
requests, binds request data to method parameters or Java objects, prints results in the server
console, and returns either plain text, a JSP page, or JSON data.

The main files developed in this experiment are:

\begin{itemize}
    \item \texttt{StudentController.java}
    \item \texttt{Student.java}
    \item \texttt{Course.java}
    \item \texttt{result.jsp}
\end{itemize}

\section{Implementation}

\subsection{Student Model}

The \texttt{Student} class is used for POJO binding and JSON response. It contains three fields:
\texttt{id}, \texttt{name}, and \texttt{age}.

\begin{lstlisting}[caption={Student.java},label={lst:student},language=Java]
package com.example.model;

public class Student {
    private Integer id;
    private String name;
    private Integer age;

    public Integer getId() {
        return id;
    }

    public void setId(Integer id) {
        this.id = id;
    }

    public String getName() {
        return name;
    }

    public void setName(String name) {
        this.name = name;
    }

    public Integer getAge() {
        return age;
    }

    public void setAge(Integer age) {
        this.age = age;
    }
}
\end{lstlisting}

\subsection{Course Model}

The \texttt{Course} class is used for JSON request binding. It contains \texttt{courseId},
\texttt{courseName}, and \texttt{credit}.

\begin{lstlisting}[caption={Course.java},label={lst:course},language=Java]
package com.example.model;

public class Course {
    private Integer courseId;
    private String courseName;
    private Integer credit;

    public Integer getCourseId() {
        return courseId;
    }

    public void setCourseId(Integer courseId) {
        this.courseId = courseId;
    }

    public String getCourseName() {
        return courseName;
    }

    public void setCourseName(String courseName) {
        this.courseName = courseName;
    }

    public Integer getCredit() {
        return credit;
    }

    public void setCredit(Integer credit) {
        this.credit = credit;
    }
}
\end{lstlisting}

\subsection{Student Controller}

The \texttt{StudentController} contains all controller methods required for this experiment,
including simple parameter binding, \texttt{@RequestParam}, \texttt{@PathVariable}, POJO binding,
list binding, JSON request binding, and JSON response.

\begin{lstlisting}[caption={StudentController.java},label={lst:controller},language=Java]
package com.example.controller;

import com.example.model.Course;
import com.example.model.Student;
import org.springframework.stereotype.Controller;
import org.springframework.ui.Model;
import org.springframework.web.bind.annotation.GetMapping;
import org.springframework.web.bind.annotation.PathVariable;
import org.springframework.web.bind.annotation.PostMapping;
import org.springframework.web.bind.annotation.RequestBody;
import org.springframework.web.bind.annotation.RequestMapping;
import org.springframework.web.bind.annotation.RequestParam;
import org.springframework.web.bind.annotation.ResponseBody;

import java.util.List;

@Controller
@RequestMapping("/student")
public class StudentController {

    @GetMapping("/info")
    @ResponseBody
    public String info(String name, int age) {
        System.out.println("Student name: " + name);
        System.out.println("Student age: " + age);
        return "Student info received";
    }

    @GetMapping("/search")
    @ResponseBody
    public String search(@RequestParam("studentName") String name) {
        System.out.println("Searching student: " + name);
        return "Search request received";
    }

    @GetMapping("/{id:\\d+}")
    @ResponseBody
    public String findById(@PathVariable("id") Integer id) {
        System.out.println("Student ID: " + id);
        return "Student ID received";
    }

    @PostMapping("/register")
    public String register(Student student, Model model) {
        model.addAttribute("student", student);
        return "result";
    }

    @PostMapping("/selectCourses")
    @ResponseBody
    public String selectCourses(@RequestParam("courses") List<String> courses) {
        System.out.println("Selected courses:");
        for (String course : courses) {
            System.out.println(course);
        }
        return "Courses received";
    }

    @PostMapping("/addCourse")
    @ResponseBody
    public String addCourse(@RequestBody Course course) {
        System.out.println("Course ID: " + course.getCourseId());
        System.out.println("Course Name: " + course.getCourseName());
        System.out.println("Credit: " + course.getCredit());
        return "Course received";
    }

    @GetMapping("/json")
    @ResponseBody
    public Student json() {
        Student student = new Student();
        student.setId(1001);
        student.setName("Tom");
        student.setAge(20);
        return student;
    }
}
\end{lstlisting}

\subsection{Result JSP Page}

The \texttt{result.jsp} page displays the result of student registration after POJO binding.

\begin{lstlisting}[caption={result.jsp},label={lst:resultjsp},language=HTML]
<%@ page contentType="text/html;charset=UTF-8" language="java" %>
<!DOCTYPE html>
<html>
<head>
    <meta charset="UTF-8">
    <title>Registration Result</title>
</head>
<body>
    <h2>Registration Successful</h2>

    <p>Student ID: ${student.id}</p>
    <p>Student Name: ${student.name}</p>
    <p>Student Age: ${student.age}</p>
</body>
</html>
\end{lstlisting}

\section{Testing and Results}

The project was started using the following Maven command:

\begin{lstlisting}[caption={Running the application}]
mvn jetty:run
\end{lstlisting}

After the server started successfully, the following tests were performed.

\subsection{Simple Parameter Binding Test}

The simple parameter binding test uses a GET request:

\begin{lstlisting}
GET http://localhost:8080/student/info?name=Tom&age=20
\end{lstlisting}

The controller receives \texttt{name} and \texttt{age} directly as method parameters and prints
the values in the server console.

\screenshot{simple-parameter-binding.png}
{Simple parameter binding test in browser}
{fig:simple-parameter-binding}

\screenshot{simple-parameter-console.png}
{Console output for simple parameter binding}
{fig:simple-parameter-console}

\subsection{POJO Data Binding Test}

The POJO data binding test uses a POST request:

\begin{lstlisting}
POST http://localhost:8080/student/register?id=1001&name=Tom&age=20
\end{lstlisting}

Spring MVC binds the request parameters to a \texttt{Student} object. The object is added to the
\texttt{Model}, and the request is forwarded to \texttt{result.jsp}.

\screenshot{pojo-binding-register.png}
{POJO data binding test with student registration result}
{fig:pojo-binding-register}

\subsection{Multiple Courses Binding Test}

The multiple courses test uses a POST request with repeated parameter names:

\begin{lstlisting}
POST http://localhost:8080/student/selectCourses
courses=Java
courses=Spring MVC
courses=MyBatis
\end{lstlisting}

Spring MVC binds the repeated \texttt{courses} parameters to a \texttt{List<String>}.

\screenshot{multiple-courses-binding.png}
{Multiple courses binding test in Postman}
{fig:multiple-courses-binding}

\screenshot{multiple-courses-console.png}
{Console output for multiple courses binding}
{fig:multiple-courses-console}

\subsection{JSON Response Test}

The JSON response test uses a GET request:

\begin{lstlisting}
GET http://localhost:8080/student/json
\end{lstlisting}

The controller returns a \texttt{Student} object with \texttt{@ResponseBody}. Spring MVC converts
the Java object into JSON.

\screenshot{json-response.png}
{JSON response test in browser}
{fig:json-response}

\section{Discussion}

The experiment shows that Spring MVC provides several convenient methods for receiving request
data. Simple request parameters can be bound directly to method arguments. When the request
parameter name and method parameter name are different, \texttt{@RequestParam} can be used.
REST-style path data can be obtained using \texttt{@PathVariable}. Complex form data can be bound
directly to a Java object, and repeated parameters can be bound to a list. JSON request and
response handling is supported through \texttt{@RequestBody} and \texttt{@ResponseBody}.

\section{Conclusion}

This experiment successfully implemented the required Student Course Registration System using
Spring MVC. The application handled simple parameter binding, POJO binding, list binding, JSON
request binding, and JSON response generation. The screenshots show successful execution of the
required tests, including simple parameter binding, POJO binding, multiple courses binding, and
JSON response.

\end{document}
